% =============================================================
%  CHAPTER 6 - REALIZATION
% =============================================================

\chapteroverview{Realization}{
  \item Introduction
  \item Hardware Environment
  \item Software Environment
  \item Adopted Architecture
  \item Database Schema
  \item Project Statistics
  \item Representative Interfaces
  \item Tests and Validation
  \item Conclusion
}

\section{Introduction}

After describing \textit{what} the platform does, this last chapter explains \textit{how} it has been built. We present the hardware and software environments, the adopted architecture, the database schema, a few representative interfaces and the validation strategy that supports the quality of the final product.

% ---------- 6.2 hardware ----------
\section{Hardware Environment}

\begin{table}[H]
\centering
\renewcommand{\arraystretch}{1.35}
\begin{tabularx}{\textwidth}{|>{\columncolor{softblue}}l|X|}
\hline
\rowcolor{primary}
\textbf{\color{white} Component} & \textbf{\color{white} Specification}\\
\hline
Processor & Intel Core i5 / i7 (10th generation or higher) \\\hline
RAM       & 8 -- 16 GB DDR4 \\\hline
Storage   & 256 -- 512 GB SSD \\\hline
Display   & 14" -- 15.6" Full HD \\\hline
Network   & Wi-Fi 5 / Wi-Fi 6, Ethernet \\\hline
\end{tabularx}
\caption{Hardware environment used during development}
\label{tab:hw}
\end{table}

% ---------- 6.3 software ----------
\section{Software Environment}

\subsection{Languages and frameworks}

\begin{figure}[H]
\centering
\begin{tikzpicture}[
  every node/.style={font=\small},
  techbox/.style={rectangle, draw=primary, fill=softblue, rounded corners=4pt,
                  minimum width=2.5cm, minimum height=0.9cm, align=center}
]
  \node[techbox] (a) at (0,0)   {HTML5};
  \node[techbox] (b) at (3,0)   {CSS3};
  \node[techbox] (c) at (6,0)   {JavaScript};
  \node[techbox] (d) at (9,0)   {React.js};
  \node[techbox] (e) at (0,-1.5){Node.js};
  \node[techbox] (f) at (3,-1.5){Express.js};
  \node[techbox] (g) at (6,-1.5){MongoDB};
  \node[techbox] (h) at (9,-1.5){JWT};
  \node[techbox] (i) at (0,-3)  {Tailwind};
  \node[techbox] (j) at (3,-3)  {Chart.js};
  \node[techbox] (k) at (6,-3)  {Multer};
  \node[techbox] (l) at (9,-3)  {Nodemailer};
\end{tikzpicture}
\caption{Technical stack of the ADVANCIA platform}
\label{fig:stack}
\end{figure}

\subsection{Tools}

\begin{table}[H]
\centering
\renewcommand{\arraystretch}{1.35}
\begin{tabularx}{\textwidth}{|>{\columncolor{softblue}}l|X|}
\hline
\rowcolor{primary}
\textbf{\color{white} Tool} & \textbf{\color{white} Usage}\\
\hline
Visual Studio Code & Main IDE for front-end and back-end development. \\\hline
Git \& GitHub      & Source-code management and collaboration. \\\hline
Postman            & Manual testing of the REST API. \\\hline
Figma              & UI/UX design and prototyping. \\\hline
Draw.io / StarUML  & UML modeling. \\\hline
Trello / Jira      & Sprint and backlog tracking. \\\hline
MongoDB Compass    & Database visualization. \\\hline
\LaTeX\ (Overleaf) & Writing of the present dissertation. \\\hline
\end{tabularx}
\caption{Tools used during the project}
\label{tab:tools}
\end{table}

% ---------- 6.4 architecture ----------
\section{Adopted Architecture}

The platform follows a classical three-tier architecture, with a clear separation between the presentation layer (React SPA), the business layer (Express REST API) and the data layer (MongoDB).

\begin{figure}[H]
\centering
\begin{tikzpicture}[
  every node/.style={font=\small},
  layer/.style={rectangle, draw=primary, fill=softblue, rounded corners=6pt,
                minimum width=12cm, minimum height=1.6cm, align=center},
  client/.style={rectangle, draw=primary, fill=white, rounded corners=4pt,
                 minimum width=2.5cm, minimum height=0.9cm, align=center},
  arr/.style={->, thick, color=primary}
]
  \node[layer] (l1) at (0,4) {\textbf{Presentation layer}\\React SPA -- Tailwind UI -- Chart.js};
  \node[layer] (l2) at (0,2) {\textbf{Business layer}\\Express REST API -- JWT -- Validation -- Mailer};
  \node[layer] (l3) at (0,0) {\textbf{Data layer}\\MongoDB collections + GridFS for files};

  \draw[arr] (l1) -- node[right]{HTTPS / JSON} (l2);
  \draw[arr] (l2) -- node[right]{Mongoose ORM}  (l3);

  \node[client, left=0.4cm of l1] (cl) {\faDesktop\ Browser};
  \draw[arr] (cl) -- (l1);
\end{tikzpicture}
\caption{Three-tier architecture of the ADVANCIA platform}
\label{fig:arch}
\end{figure}

\subsection{Deployment view}

\begin{figure}[H]
\centering
\begin{tikzpicture}[
  every node/.style={font=\small},
  srv/.style={rectangle, draw=primary, fill=softblue, rounded corners=4pt,
              minimum width=3cm, minimum height=1cm, align=center},
  arr/.style={->, thick, color=primary}
]
  \node[srv] (browser) at (0,0)   {\faDesktop\ Browser};
  \node[srv] (front)   at (4,0)   {\faServer\ Front Server\\(Vite / Nginx)};
  \node[srv] (api)     at (8,0)   {\faCogs\ API Server\\(Node + Express)};
  \node[srv] (db)      at (12,0)  {\faDatabase\ MongoDB};
  \node[srv] (mail)    at (8,-1.8){\faEnvelope\ SMTP / Mail};

  \draw[arr] (browser) -- (front);
  \draw[arr] (front)   -- (api);
  \draw[arr] (api)     -- (db);
  \draw[arr] (api)     -- (mail);
\end{tikzpicture}
\caption{Simplified deployment diagram}
\label{fig:deploy}
\end{figure}

% ---------- 6.5 DB ----------
\section{Database Schema}

\begin{figure}[H]
\centering
\resizebox{\textwidth}{!}{%
\begin{tikzpicture}[
  every node/.style={font=\small},
  ent/.style={rectangle split, rectangle split parts=2,
              rectangle split part fill={primary, white},
              draw=primary, text=white, rectangle split draw splits=true, align=left, inner sep=4pt}
]
  \node[ent] (u) at (0,0) {%
    \nodepart{one}\bfseries User
    \nodepart[text=black]{two}
      \underline{\_id}\\ email\\ password\\ firstName\\ lastName\\ photo\\ avatar\\ role
  };
  \node[ent, right=2cm of u] (t) {%
    \nodepart{one}\bfseries Training
    \nodepart[text=black]{two}
      \underline{\_id}\\ title\\ description\\ duration\\ price\\ level\\ categoryId
  };
  \node[ent, right=2cm of t] (c) {%
    \nodepart{one}\bfseries Course
    \nodepart[text=black]{two}
      \underline{\_id}\\ trainingId\\ title\\ content\\ resources\\ order
  };
  \node[ent, below=2cm of u] (r) {%
    \nodepart{one}\bfseries Reservation
    \nodepart[text=black]{two}
      \underline{\_id}\\ userId\\ trainingId\\ status\\ createdAt
  };
  \node[ent, below=2cm of t] (n) {%
    \nodepart{one}\bfseries Notification
    \nodepart[text=black]{two}
      \underline{\_id}\\ recipientId\\ type\\ message\\ read
  };
  \node[ent, below=2cm of c] (cat) {%
    \nodepart{one}\bfseries Category
    \nodepart[text=black]{two}
      \underline{\_id}\\ name\\ description
  };

  \draw[->, thick, primary] (u)   -- node[above,font=\scriptsize]{1..*}(r);
  \draw[->, thick, primary] (t)   -- node[above,font=\scriptsize]{1..*}(r);
  \draw[->, thick, primary] (t)   -- node[above,font=\scriptsize]{1..*}(c);
  \draw[->, thick, primary] (cat) -- node[right,font=\scriptsize]{1..*}(t);
  \draw[->, thick, primary] (u)   -- node[above,font=\scriptsize]{*}(n);
\end{tikzpicture}}
\caption{Entity-relationship schema of the ADVANCIA database}
\label{fig:erd}
\end{figure}

% ---------- 6.6 stats ----------
\section{Project Statistics}

\begin{table}[H]
\centering
\renewcommand{\arraystretch}{1.35}
\begin{tabularx}{\textwidth}{|>{\columncolor{softblue}}l|X|}
\hline
\rowcolor{primary}
\textbf{\color{white} Indicator} & \textbf{\color{white} Value}\\
\hline
Number of user stories      & 18 \\\hline
Number of releases          & 3 \\\hline
Number of sprints           & 6 \\\hline
Total story points          & 95 \\\hline
Front-end pages             & $\sim$ 22 \\\hline
REST API endpoints          & $\sim$ 45 \\\hline
Database collections        & 7 \\\hline
Manual test cases           & $\sim$ 60 \\\hline
\end{tabularx}
\caption{Quantitative indicators of the project}
\label{tab:stats}
\end{table}

\subsection{Code distribution}

\begin{figure}[H]
\centering
\begin{tikzpicture}
\begin{axis}[
  width=11cm, height=5.5cm,
  ybar,
  bar width=22pt,
  ylabel={Lines of code (\%)},
  symbolic x coords={Front-end, Back-end, Database, Tests, Config},
  xtick=data,
  ymin=0, ymax=55,
  nodes near coords,
  nodes near coords style={font=\small\bfseries, color=primary},
  enlarge x limits=0.15,
  every axis label/.append style={font=\small}
]
  \addplot[fill=primary,    draw=primary]    coordinates {(Front-end,42)};
  \addplot[fill=secondary,  draw=secondary]  coordinates {(Back-end,38)};
  \addplot[fill=accent,     draw=accent]     coordinates {(Database,7)};
  \addplot[fill=success,    draw=success]    coordinates {(Tests,9)};
  \addplot[fill=danger,     draw=danger]     coordinates {(Config,4)};
\end{axis}
\end{tikzpicture}
\caption{Approximate distribution of the code base}
\label{fig:codedist}
\end{figure}

\subsection{Sprint velocity}

\begin{figure}[H]
\centering
\begin{tikzpicture}
\begin{axis}[
  width=12cm, height=5.5cm,
  ymajorgrids, grid style=dashed,
  ylabel={Story points completed},
  symbolic x coords={S1, S2, S3, S4, S5, S6},
  xtick=data,
  ymin=0, ymax=22,
  mark size=2.5pt,
  nodes near coords,
  nodes near coords style={font=\small\bfseries, color=primary}
]
  \addplot[mark=*, color=primary, thick] coordinates {
    (S1,8) (S2,12) (S3,14) (S4,18) (S5,17) (S6,20)
  };
\end{axis}
\end{tikzpicture}
\caption{Velocity per sprint (story points completed)}
\label{fig:velocity}
\end{figure}

% ---------- 6.7 screens ----------
\section{Representative Interfaces}

This section illustrates the main interfaces produced. Real screenshots will replace the placeholders in the final printed version of the dissertation.

\begin{figure}[H]
\centering
\begin{tikzpicture}[
  every node/.style={font=\small},
  scr/.style={rectangle, draw=primary, fill=softblue, rounded corners=6pt,
              minimum width=5.5cm, minimum height=3cm, align=center,
              font=\small\bfseries\color{primary}}
]
  \node[scr] (a) at (0,0)   {Sign-up / Login\\page};
  \node[scr] (b) at (6.2,0) {User\\dashboard};
  \node[scr] (c) at (0,-3.5){Catalog of\\trainings};
  \node[scr] (d) at (6.2,-3.5){Reservation\\workflow};
  \node[scr] (e) at (0,-7)  {Smart\\admin dashboard};
  \node[scr] (f) at (6.2,-7){Chatbot\\assistant};
\end{tikzpicture}
\caption{Mosaic of the main interfaces of the platform}
\label{fig:mosaic}
\end{figure}

% ---------- 6.8 tests ----------
\section{Tests and Validation}

A pyramid testing strategy was followed during the project : unit tests at the bottom (most numerous, fastest), integration tests in the middle, and end-to-end tests at the top (less numerous, slower but representative of the user journey).

\begin{figure}[H]
\centering
\begin{tikzpicture}
  \fill[primary]   (-3,0) -- (3,0)  -- (2,2) -- (-2,2) -- cycle;
  \fill[secondary] (-2,2) -- (2,2)  -- (1,4) -- (-1,4) -- cycle;
  \fill[accent]    (-1,4) -- (1,4)  -- (0,5.5) -- cycle;

  \node[white, font=\bfseries] at (0,1)   {Unit tests};
  \node[white, font=\bfseries] at (0,3)   {Integration tests};
  \node[white, font=\bfseries] at (0,4.7) {E2E tests};
\end{tikzpicture}
\caption{Testing pyramid adopted}
\label{fig:pyramid}
\end{figure}

\begin{table}[H]
\centering
\renewcommand{\arraystretch}{1.35}
\begin{tabularx}{\textwidth}{|>{\columncolor{softblue}}l|c|c|X|}
\hline
\rowcolor{primary}
\textbf{\color{white} Test} & \textbf{\color{white} Cases} & \textbf{\color{white} Pass rate} & \textbf{\color{white} Coverage}\\
\hline
Unit (services)      & 38 & 100 \% & Auth, reservation, dashboard \\\hline
Integration (API)    & 16 & 100 \% & Critical endpoints \\\hline
E2E (UI)             & 8  & 100 \% & Sign-up, reserve, accept, export \\\hline
\end{tabularx}
\caption{Synthesis of the validation tests}
\label{tab:tests}
\end{table}

\section{Conclusion}

This realization chapter has materialized the project on a technical level : a modern stack, a clean three-tier architecture, a documented database, measurable indicators and a structured testing strategy. The platform is now ready for delivery and for the perspectives of evolution detailed in the general conclusion.
